% GENERATED FILE. Edit canonical lessons, then run npm run generate:books.
% canonical-source-hash: fnv1a64:7feb2de4aaf2d767
% canonical-lessons: TA-C360-niyayam, TA-C360-niti, TA-C360-aram, TA-C360-pali, TA-C360-cakku

\chapter{Right and Wrong}
\label{ch:ta-right-and-wrong}
\hlchaptermodality{360}

\begin{chapteropening}
\textbf{By the end of this chapter:} \emph{I can talk about fairness, justice, blame and excuses.}

Complete the last lesson of chapter 360: I can talk about fairness, justice, blame and excuses.
\end{chapteropening}

\section[\texorpdfstring{\ta{நியாயம்} (niyāyam)}{நியாயம் (niyāyam)}]{\ta{நியாயம்} (niyāyam) — fairness; what is right}

\label{lesson:TA-C360-niyayam}



\begin{quote}
Before the new one: say the Tamil for uncertainty, then the Tamil for to reject.
\end{quote}

\subsection*{You'll want to know: \ta{நியாயம்}}
\textbf{\ta{நியாயம்}} — \emph{niyāyam} — \textquotedblleft{}fairness; what is right\textquotedblright{}.

\begin{quote}
\ta{அவருக்கு} \ta{நடந்தது} \ta{நியாயம்} \ta{இல்லை}.
\end{quote}

\begin{quote}
\emph{avarukku naṭantatu niyāyam illai.}
\end{quote}

\begin{quote}
What happened to him is not fair.
\end{quote}

\ta{நியாயமான}, fair: \ta{நியாயமான} \ta{விலை}, a fair price.

\begin{cousinweb}
\textbf{\ta{நியாயம்}} (\emph{niyāyam}) is borrowed from Sanskrit \emph{nyāya}, \textquotedblleft{}rule, method\textquotedblright{} (CDIAL 7614); Pali has \emph{ñāya}, \textquotedblleft{}justice\textquotedblright{}. Tamil breaks the cluster ny with a vowel: niyāyam.
\end{cousinweb}

\begin{grammarlens}[title={What you've built}]
One more word for giving reasons and opinions.
\end{grammarlens}

\subsection*{Your turn}
\begin{itemize}
\raggedright
  \item \emph{Say it:} \emph{niyāyam}
  \item \emph{Say it:} \emph{niyāyam}, once more
  \item \emph{Say it:} \emph{avarukku naṭantatu niyāyam illai.} — What happened to him is not fair.
  \item \emph{Recall:} say the Tamil for uncertainty, then the Tamil for to reject, then say \emph{niyāyam} again
\end{itemize}

\begin{tcolorbox}[breakable,colback=teal!4,colframe=teal!35!black,title={Before you move on}]
What does \ta{நியாயம்} mean? (\textquotedblleft{}Fairness; what is right\textquotedblright{}.) Say it once more.
\end{tcolorbox}
\section[\texorpdfstring{\ta{நீதி} (nīti)}{நீதி (nīti)}]{\ta{நீதி} (nīti) — justice}

\label{lesson:TA-C360-niti}



\begin{quote}
Before the new one: say the Tamil for to reject, then the Tamil for fairness.
\end{quote}

\subsection*{You'll want to know: \ta{நீதி}}
\textbf{\ta{நீதி}} — \emph{nīti} — \textquotedblleft{}justice\textquotedblright{}.

\begin{quote}
\ta{பாதிக்கப்பட்ட} \ta{மக்களுக்கு} \ta{நீதி} \ta{கிடைக்க} \ta{வேண்டும்}.
\end{quote}

\begin{quote}
\emph{pātikkappaṭṭa makkaḷukku nīti kiṭaikka vēṇṭum.}
\end{quote}

\begin{quote}
The people who were harmed must get justice.
\end{quote}

\ta{நீதி} \ta{கேள்}, to demand justice. \ta{நீதிபதி}, a judge, is \ta{நீதி} + \ta{பதி}, master.

\begin{cousinweb}
\textbf{\ta{நீதி}} (\emph{nīti}) is borrowed from Sanskrit \emph{nīti}, \textquotedblleft{}conduct, guidance\textquotedblright{} (CDIAL 7546), from the root \emph{nī}, \textquotedblleft{}to lead\textquotedblright{}; Sinhala \emph{niya} means \textquotedblleft{}law\textquotedblright{}.
\end{cousinweb}

\begin{grammarlens}[title={What you've built}]
One more word for giving reasons and opinions.
\end{grammarlens}

\subsection*{Your turn}
\begin{itemize}
\raggedright
  \item \emph{Say it:} \emph{nīti}
  \item \emph{Say it:} \emph{nīti}, once more
  \item \emph{Say it:} \emph{pātikkappaṭṭa makkaḷukku nīti kiṭaikka vēṇṭum.} — The people who were harmed must get justice.
  \item \emph{Recall:} say the Tamil for to reject, then the Tamil for fairness, then say \emph{nīti} again
\end{itemize}

\begin{tcolorbox}[breakable,colback=teal!4,colframe=teal!35!black,title={Before you move on}]
What does \ta{நீதி} mean? (\textquotedblleft{}Justice\textquotedblright{}.) Say it once more.
\end{tcolorbox}
\section[\texorpdfstring{\ta{அறம்} (aṟam)}{அறம் (aṟam)}]{\ta{அறம்} (aṟam) — ethics; virtue}

\label{lesson:TA-C360-aram}



\begin{quote}
Before the new one: say the Tamil for fairness, then the Tamil for justice.
\end{quote}

\subsection*{You'll want to know: \ta{அறம்}}
\textbf{\ta{அறம்}} — \emph{aṟam} — \textquotedblleft{}ethics; virtue\textquotedblright{}.

\begin{quote}
\ta{அறம்} \ta{இல்லாத} \ta{வாழ்க்கை} \ta{பயனற்றது}.
\end{quote}

\begin{quote}
\emph{aṟam illāta vāḻkkai payaṉaṟṟatu.}
\end{quote}

\begin{quote}
A life without virtue is worthless.
\end{quote}

An old Tamil word; \ta{அறவழி}, the path of virtue. \ta{நீதி} is justice between people, \ta{அறம்} the good in a person.

\begin{cousinweb}
\textbf{\ta{அறம்}} (\emph{aṟam}) is native Dravidian: DEDR 311 lists Tamil \emph{aṟam}, \textquotedblleft{}moral duty, virtue\textquotedblright{}, and \emph{aṟavaṉ}, \textquotedblleft{}one who is virtuous\textquotedblright{}.
\end{cousinweb}

\begin{grammarlens}[title={What you've built}]
One more word for giving reasons and opinions.
\end{grammarlens}

\subsection*{Your turn}
\begin{itemize}
\raggedright
  \item \emph{Say it:} \emph{aṟam}
  \item \emph{Say it:} \emph{aṟam}, once more
  \item \emph{Say it:} \emph{aṟam illāta vāḻkkai payaṉaṟṟatu.} — A life without virtue is worthless.
  \item \emph{Recall:} say the Tamil for fairness, then the Tamil for justice, then say \emph{aṟam} again
\end{itemize}

\begin{tcolorbox}[breakable,colback=teal!4,colframe=teal!35!black,title={Before you move on}]
What does \ta{அறம்} mean? (\textquotedblleft{}Ethics; virtue\textquotedblright{}.) Say it once more.
\end{tcolorbox}
\section[\texorpdfstring{\ta{பழி} (paḻi)}{பழி (paḻi)}]{\ta{பழி} (paḻi) — blame}

\label{lesson:TA-C360-pali}



\begin{quote}
Before the new one: say the Tamil for justice, then the Tamil for ethics.
\end{quote}

\subsection*{You'll want to know: \ta{பழி}}
\textbf{\ta{பழி}} — \emph{paḻi} — \textquotedblleft{}blame\textquotedblright{}.

\begin{quote}
\ta{அவன்} \ta{தன்} \ta{தவறுக்கு} \ta{மற்றவர்கள்} \ta{மீது} \ta{பழி} \ta{போட்டான்}.
\end{quote}

\begin{quote}
\emph{avaṉ taṉ tavaṟukku maṟṟavarkaḷ mītu paḻi pōṭṭāṉ.}
\end{quote}

\begin{quote}
He blamed others for his mistake.
\end{quote}

X \ta{மீது} \ta{பழி} \ta{போடு}, to put the blame on X.

\begin{cousinweb}
\textbf{\ta{பழி}} (\emph{paḻi}) is native Dravidian: DEDR 4002 lists Tamil \emph{paḻi}, \textquotedblleft{}to blame, censure; blame, fault\textquotedblright{}, with Malayalam \emph{paḻi}, \textquotedblleft{}fault, false accusation\textquotedblright{}.
\end{cousinweb}

\begin{grammarlens}[title={What you've built}]
One more word for giving reasons and opinions.
\end{grammarlens}

\subsection*{Your turn}
\begin{itemize}
\raggedright
  \item \emph{Say it:} \emph{paḻi}
  \item \emph{Say it:} \emph{paḻi}, once more
  \item \emph{Say it:} \emph{avaṉ taṉ tavaṟukku maṟṟavarkaḷ mītu paḻi pōṭṭāṉ.} — He blamed others for his mistake.
  \item \emph{Recall:} say the Tamil for justice, then the Tamil for ethics, then say \emph{paḻi} again
\end{itemize}

\begin{tcolorbox}[breakable,colback=teal!4,colframe=teal!35!black,title={Before you move on}]
What does \ta{பழி} mean? (\textquotedblleft{}Blame\textquotedblright{}.) Say it once more.
\end{tcolorbox}
\section[\texorpdfstring{\ta{சாக்கு} (cākku)}{சாக்கு (cākku)}]{\ta{சாக்கு} (cākku) — an excuse; a pretext}

\label{lesson:TA-C360-cakku}



\begin{quote}
Before the new one: say the Tamil for ethics, then the Tamil for blame.
\end{quote}

\subsection*{You'll want to know: \ta{சாக்கு}}
\textbf{\ta{சாக்கு}} — \emph{cākku} — \textquotedblleft{}an excuse; a pretext\textquotedblright{}.

\begin{quote}
\ta{தாமதமாக} \ta{வந்ததற்கு} \ta{அவன்} \ta{ஒரு} \ta{சாக்கு} \ta{சொன்னான்}.
\end{quote}

\begin{quote}
\emph{tāmatamāka vantataṟku avaṉ oru cākku coṉṉāṉ.}
\end{quote}

\begin{quote}
He made an excuse for coming late.
\end{quote}

\ta{சாக்குப்போக்கு}, excuses, an echo pair. \ta{சாக்கு}, a sack, sounds the same but is another word.

\begin{cousinweb}
\textbf{\ta{சாக்கு}} (\emph{cākku}) is native Dravidian: DEDR 2428 lists Tamil \emph{cākku}, \textquotedblleft{}excuse, pretext\textquotedblright{}, with Kannada \emph{sāku} and Telugu \emph{sāku}, \textquotedblleft{}pretence, plea\textquotedblright{}.
\end{cousinweb}

\begin{grammarlens}[title={What you've built}]
One more word for giving reasons and opinions.
\end{grammarlens}

\subsection*{Your turn}
\begin{itemize}
\raggedright
  \item \emph{Say it:} \emph{cākku}
  \item \emph{Say it:} \emph{cākku}, once more
  \item \emph{Say it:} \emph{tāmatamāka vantataṟku avaṉ oru cākku coṉṉāṉ.} — He made an excuse for coming late.
  \item \emph{Recall:} say the Tamil for ethics, then the Tamil for blame, then say \emph{cākku} again
\end{itemize}

\begin{tcolorbox}[breakable,colback=teal!4,colframe=teal!35!black,title={Before you move on}]
What does \ta{சாக்கு} mean? (\textquotedblleft{}An excuse; a pretext\textquotedblright{}.) Say it once more.
\end{tcolorbox}
